
\subsection*{4.4.}
On
appelle parfois l'entier (4.3.4) le \emph{nombre d'intersection} des \(\underline{O}_X\) modules inversibles \(\underline{L}_1,\ldots,\underline{L}_d\) avec \(\underline{F}\), et on le dénote par \((\underline{L}_1,\ldots,\underline{L}_d;\underline{F})\) (cf. [3]). Comme conséquence de l'expression (4.3.4) de ce nombre on trouve aussitôt les propriétés suivantes :

\begin{enumerate}[label=\arabic*)]
\item L'entier (4.3.4) est une fonction symétrique et \(d\)-linéaire dans les variables \(L_1,\ldots,L_d\in\Pic(X)\), nulle si \(\dim(\Supp(\underline{F}))<d\), et il est «~additif~» en \(\underline{F}\) i.e. additif en \(\cl_\bullet(\underline F)=F\).

\item Utilisant 2.1.2, on en conclut immédiatement (avec les notations de 1.1.2) :
\begin{equation}\tag{4.4.1}
 (\underline{L}_1,\ldots,\underline{L}_d;\underline{F})
 =
 \sum_{x\in X(d)}
 \operatorname{long}_{\underline{O}_{,x}}(\underline{F}_x)
 (\underline{L}_1,\ldots,\underline{L}_d;\underline{O}_{Z(x)}).
\end{equation}

\item Soit \(Y\) un sous-schéma fermé de \(X\) et \(\underline{G}\) un \(\underline{O}_Y\)-module cohérent tel que \(\underline{F}=f_*(\underline{G})\), où \(f:Y\to X\) est l'immersion canonique. Alors on trouve
\[
 (f^*\underline{L}_1,\ldots,f^*\underline{L}_d;\underline{G})=(\underline{L}_1,\ldots,\underline{L}_d;\underline{F}).
\]
De façon plus générale, la formule de projection donne en effet
\begin{equation}\tag{4.4.2}
 \chi(f^*(x)\cdot y)=\chi\bigl(f_*(f^*(x)\cdot y)\bigr)=\chi(x\cdot f_*(y))
 \qquad (x\in K^{\bullet}(X),\ y\in K_\bullet(Y))
\end{equation}
pour tout \(k\)-morphisme \(f:Y\to X\) de \(k\)-schémas propres.

\item Supposons que dans le morphisme \(f:Y\to X\) ci-dessus, \(X\) et \(Y\) sont irréductibles de dimension \(\leq d\). Si \(f\) est de degré \(n\), alors on a
